\paragraph{$^{$FOOTNOTE_NUMBER$}$Averaged quantities}\label{footnote:$FOOTNOTE_NUMBER$} are calculated by accumulating the quantity over
the 10Myr/100Myr that precedes the writing of a snapshot, and then normalising. For example for SFR we start a clock precisely 10Myr before a snapshot
dump, accumulate SFR * dt at each step during that window, and then divide by 10Myr at the point of writing.

For \verb+AveragedStarFormationRate+ we sum the averaged SFRs of both the gas and
the star particles. When a gas particle is converted into a star particle the value
it has accumulated so far is copied to the star particle, and accumulation then
stops. Star particles that formed before the start of the averaging window have a
value of zero. The value of each star particle is normalised by the length of
the averaging window (not the time since the star formed), so summing the contribution of the gas and star particles
gives the total star formation rate averaged over the window.

